\subsection{Environment Configuration}

All evaluations use a fixed RetailBench configuration designed to capture the complexity of sustained retail operations. The environment comprises 96 products across 20 categories, with five candidate suppliers per product whose price--quality profiles evolve over time. Customer reviews and daily news events introduce additional operational feedback. Agents operate with an initial budget of 30{,}000, daily rent of 600, an inventory capacity of 15{,}000 units, and a shelf capacity of 40 products. Each rollout spans up to 180 days and terminates upon bankruptcy,
defined as five consecutive days of negative funds. Complete environment settings are provided in Appendix~\ref{appendix:env_setting}.


\subsection{Models and Scaffolds}
\label{sec:models_frameworks}

We evaluate LLM agents using three prompting-based scaffolds: ReAct~\citep{yao2023react}, Reflection~\citep{reflection}, and Plan-and-Act~\citep{erdogan2025planandactimprovingplanningagents}. These scaffolds represent common agent interaction patterns and operate through the same RetailBench tool interface; Appendix~\ref{appendix:prompts} lists their prompts.

Our evaluation covers seven LLMs and 19 model--scaffold configurations: six models under all three scaffolds and GPT-5.5 under ReAct only. For models with complete scaffold coverage, we report means and sample standard deviations across scaffolds, reflecting scaffold sensitivity rather than seed-level uncertainty. 


\subsection{Reference Policies}
\label{sec:references}

We compare LLM agents against two non-LLM reference policies that differ in information access: OR-Composite uses only agent-visible information, whereas the Privileged Heuristic reads hidden simulator state.

\paragraph{OR-Composite.}
OR-Composite is a classical operations-research reference restricted to agent-visible information. It combines textbook rules for four decisions: a fixed assortment of at most two SKUs per category ranked by visible historical contribution; quality-first supplier selection~\citep{dickson1966vendor,weber1991vendor} that ranks vendors by an upper confidence bound on observed return rates~\citep{brown2001interval}; daily-review $(s,S)$ replenishment capped by shelf life~\citep{silver2017inventory,nahmias1982perishable}; and weekly constant-elasticity markup pricing over return-adjusted costs~\citep{phillips2021pricing}. It never reads latent supplier quality, hidden demand, news effects, or future outcomes. Appendix~\ref{appendix:or_composite} gives the full specification, parameters, and information-access audit.

\paragraph{Privileged Heuristic.}
The Privileged Heuristic is a hand-specified, rule-based reference with access to hidden simulator state. It reads structured simulator state that agents cannot observe, including supplier quality scores and news impact values: it orders each product from its highest-quality supplier, sets prices that maximize historical profit, and scales replenishment to news-adjusted demand targets. It is neither a fair baseline for LLM agents nor an upper bound, because it follows a fixed hand-specified rule and faces none of the agents' prompt, tool-selection, or memory constraints. Appendix~\ref{appendix:handcrafted_policy} gives its algorithm and parameter settings.

% Table 2 is defined here so that it floats to the top of page 6.
\begin{table*}[t]
\centering
\footnotesize
\setlength{\tabcolsep}{3.0pt}
\renewcommand{\arraystretch}{1.08}
\resizebox{\textwidth}{!}{%
\begin{tabular}{lrrrrrrr}
\toprule
Entry & Days $\uparrow$ & Net worth $\uparrow$ & Sales $\uparrow$ & Return $\downarrow$ & Expired $\downarrow$ & Stockout-day rate $\downarrow$ & Tokens per Day $\downarrow$ \\
\midrule
DeepSeek-V4-Pro & \secondval{139.3 $\pm$ 43.2} & 2,258.10 $\pm$ 6,810.42 & \secondval{110,980 $\pm$ 68,622} & \secondval{0.094 $\pm$ 0.007} & 0.039 $\pm$ 0.022 & 0.659 $\pm$ 0.111 & 450,693 $\pm$ 62,795 \\
GLM-5.1 & 55.0 $\pm$ 5.6 & -113.13 $\pm$ 4,113.31 & 4,842 $\pm$ 3,555 & \secondval{0.094 $\pm$ 0.053} & 0.013 $\pm$ 0.015 & \bestval{0.300 $\pm$ 0.148} & 1,016,126 $\pm$ 157,931 \\
Grok-4.3 & 40.3 $\pm$ 24.8 & \secondval{8,201.74 $\pm$ 12,806.30} & 10,861 $\pm$ 10,275 & 0.095 $\pm$ 0.038 & 0.077 $\pm$ 0.088 & \secondval{0.395 $\pm$ 0.036} & \secondval{259,191 $\pm$ 171,982} \\
Kimi-K2.6 & 95.0 $\pm$ 32.4 & 1,566.00 $\pm$ 4,005.87 & 55,317 $\pm$ 28,258 & 0.109 $\pm$ 0.016 & 0.014 $\pm$ 0.023 & 0.533 $\pm$ 0.103 & 684,782 $\pm$ 295,717 \\
MiniMax-M2.5 & 63.3 $\pm$ 8.7 & 602.42 $\pm$ 1,104.83 & 20,908 $\pm$ 3,587 & 0.105 $\pm$ 0.002 & \bestval{0.005 $\pm$ 0.003} & 0.731 $\pm$ 0.081 & 802,006 $\pm$ 419,031 \\
Qwen3.5-397B-A17B & 57.0 $\pm$ 12.1 & 3,709.77 $\pm$ 2,473.06 & 32,857 $\pm$ 2,774 & 0.106 $\pm$ 0.010 & 0.069 $\pm$ 0.036 & 0.645 $\pm$ 0.057 & 886,581 $\pm$ 179,962 \\
GPT-5.5 & \bestval{180.0 $\pm$ --} & \bestval{24,350.98 $\pm$ --} & \bestval{136,405 $\pm$ --} & \bestval{0.065 $\pm$ --} & \secondval{0.011 $\pm$ --} & 0.761 $\pm$ -- & \bestval{176,260 $\pm$ --} \\
\midrule
OR-Composite (classical OR) & 180.0 & 80,801.21 & 202,557 & 0.047 & 0.013 & 0.417 & \\
\midrule
Oracle Policy & 180.0 & 124,491.15 & 263,310 & 0.020 & 0.024 & 0.678 & \\
\bottomrule
\end{tabular}
}
\caption{Model-level outcome summary.
Results are reported as mean $\pm$ sample standard deviation
across three scaffolds, except GPT-5.5 (ReAct only).
LLM agents substantially underperform both references,
with the best LLM achieving only 30.1\% of OR-Composite's
final net worth.
Bold and underline indicate the best and second-best LLM results.
OR-Composite uses only agent-visible information
(Appendix~\ref{appendix:or_composite}),
while Oracle is a privileged reference
(Appendix~\ref{appendix:oracle}).}
\label{tab:aggregate_main_results}
\vspace{-3mm}
\end{table*}


\subsection{Outcome and Resource Metrics}

We evaluate long-horizon operational performance through three categories of business outcomes and two categories of resource diagnostics. Arrows indicate the preferred direction for outcome metrics.

\begin{itemize}[leftmargin=*]

    \item \textbf{Operational viability.}
    \textbf{Final Net Worth} ($\uparrow$) measures terminal asset value, including cash, on-hand inventory, and pending-order value. \textbf{Survival Days} ($\uparrow$) measures the duration of continued store operation.

    \item \textbf{Sales effectiveness.}
    \textbf{Final Total Sales} ($\uparrow$) measures cumulative units sold; and \textbf{Return Ratio} ($\downarrow$) measures the fraction of sold units that are returned.

    \item \textbf{Inventory control.}
    \textbf{Stockout Days} ($\downarrow$) counts days on which at least one product has insufficient inventory, while \textbf{Expired Ratio} ($\downarrow$) measures the fraction of expired units among all units sold or expired.

    \item \textbf{Token use.}
    \textbf{Tokens per Day} measures average token consumption as a resource-efficiency diagnostic.

\end{itemize}


\subsection{Four-Stage Diagnosis}
\label{sec:four_stage_setup}

Outcome metrics quantify long-horizon operational performance but provide limited insight into where agent decision-making breaks down. To investigate these limitations, we introduce a four-stage trajectory analysis covering product monitoring, evidence collection, decision-making, and anomaly response. Each stage is evaluated using an auditable, trace-based metric that captures a distinct aspect of operational behavior.

\begin{itemize}[leftmargin=*]

    \item \textbf{Stage 1: Product Monitoring.}
    Measures product-level inspection coverage as the fraction of products sold on day $t-1$ that receive a successful product-specific query on day $t$. Store-wide listings are excluded.

    \item \textbf{Stage 2: Evidence Collection.}
    Measures evidence completeness as the fraction of executed orders and price updates for which all required evidence dimensions were observed within their category-specific validity windows.

    \item \textbf{Stage 3: Decision-Making.}
    Measures decision-making ability in order sizing and supplier selection by averaging two shares computed from simulator state on the action day. The order-quantity share is the fraction of order lines whose post-arrival days of cover fall between the selected supplier’s maximum quoted lead time and the product’s shelf life. Days of cover are calculated by dividing the sum of inventory position and ordered quantity by expected daily demand at the agent’s current price. The supplier-quality share is the fraction of order lines assigned to the highest-quality supplier among the five candidates. These evaluator-side measures characterize ordering decisions but do not establish overall decision optimality, which also depends on pricing.
 
    
    \item \textbf{Stage 4: Anomaly Response.}
    Measures operational response coverage as the fraction of stockout and expiration events followed by an executed order or price adjustment for the affected product within three days. Events whose three-day windows extend beyond the last completed day are right-censored and excluded.

\end{itemize}

The four metrics are evaluated independently and Appendix~\ref{appendix:stage_metrics} provides the complete metric definitions, denominators, censoring rules, and finer per-stage analyses.